\subsection{Lee--Yang 分支集的超立方体谱、Haar 极限与 Ihara 闭式}

在命题 \ref{prop:killo-leyang-branchset-hypercube-families} 与定理 \ref{thm:derived-leyang-profinite-rigidity} 已把 Lee--Yang 分支集塔压缩为五个互不相交的 dyadic 超立方体族之后，还可把有限层图谱、极限计数分布以及 Ihara \(\zeta\) 函数继续写成完全显式的封闭式。

\leanverified{paper\_derived\_leyang\_branchset\_adjacency\_spectrum\_heat\_trace}
\begin{theorem}[Lee--Yang 分支集的超立方体邻接谱与热迹闭式]\label{thm:derived-leyang-branchset-adjacency-spectrum-heat-trace}
沿用命题 \ref{prop:killo-leyang-branchset-hypercube-families} 的记号，并把每个分量
\(
\Pi_{i,n}=P_{i,n}+E_{\min}[2^n]
\)
视为由 \(E_{\min}[2^n]\) 的标准生成元所诱导的 Cayley 图。则有限层分支集的自然邻接图满足
\[
R(y_n)=\bigsqcup_{i=0}^{4}\Pi_{i,n}
\cong
\bigsqcup_{i=0}^{4}Q_{2n}.
\]
因此其邻接算子 \(A_{R(y_n)}\) 的谱为
\[
\operatorname{Spec}(A_{R(y_n)})
=
\bigl\{\,2n-2j:\ 0\le j\le 2n\,\bigr\},
\qquad
\operatorname{mult}(2n-2j)=5\binom{2n}{j}.
\]
并且对任意 \(t\in\CC\)，热迹满足闭式
\[
\Tr\!\bigl(e^{tA_{R(y_n)}}\bigr)
=
5\sum_{j=0}^{2n}\binom{2n}{j}e^{t(2n-2j)}
=
5(2\cosh t)^{2n}.
\]
\end{theorem}

\begin{proof}
由命题 \ref{prop:killo-leyang-branchset-hypercube-families}，每个 \(\Pi_{i,n}\) 都与 \(E_{\min}[2^n]\) 的 Cayley 图同构；再由定理 \ref{thm:killo-leyang-2power-torsion-cayley-hypercube}，该图即 \(Q_{2n}\)。五个分支族互不相交，故 \(R(y_n)\) 正是五个 \(Q_{2n}\) 的不交并。

对 \(d\)-维超立方体 \(Q_d\)，其邻接算子可写为
\[
A_{Q_d}
=
\sum_{\ell=1}^{d}
I^{\otimes(\ell-1)}
\otimes
\begin{pmatrix}
0&1\\
1&0
\end{pmatrix}
\otimes
I^{\otimes(d-\ell)}.
\]
单坐标交换矩阵的特征值为 \(\pm1\)。因此任一张量本征向量对应的总特征值为
\(
d-2j
\),
其中 \(j\) 是取到 \(-1\) 的坐标数；其重数恰为 \(\binom{d}{j}\)。取 \(d=2n\) 并乘上五个不交并分量，即得
\[
\operatorname{mult}(2n-2j)=5\binom{2n}{j}.
\]

最后，由谱分解可得
\[
\Tr\!\bigl(e^{tA_{R(y_n)}}\bigr)
=
5\sum_{j=0}^{2n}\binom{2n}{j}e^{t(2n-2j)}
=
5e^{2nt}\sum_{j=0}^{2n}\binom{2n}{j}e^{-2tj}
=
5(e^t+e^{-t})^{2n}
=
5(2\cosh t)^{2n}.
\]
证毕。
\end{proof}

\begin{theorem}[Lee--Yang 分支集塔的有限层精确 cylinder 等分与五份 Haar 极限]\label{thm:derived-leyang-branchset-haar-limit}
沿用定理 \ref{thm:derived-leyang-profinite-rigidity} 与推论 \ref{cor:xi-terminal-zm-leyang-branchpoint-tate-torsor-decomposition} 的设定。选定每个分支族的一条相容提升 \(\mathbf P_i\in\Pi_{i,\infty}\) 及 \(T_2(E_{\min})\cong\ZZ_2^2\) 的一组 \(\ZZ_2\)-基后，有
\[
\mathfrak X_{\mathrm{LY}}
:=
\varprojlim_n \mathcal R_n
\cong
\bigsqcup_{i=0}^{4}\bigl(\mathbf P_i+\ZZ_2^2\bigr).
\]
进一步，对每个 \(n\) 以长度 \(n\) 的二进地址后接零尾的方式，把 \(\mathcal R_n\) 嵌入 \(\mathfrak X_{\mathrm{LY}}\)，并记 \(\mu_n\) 为该嵌入下归一化计数测度的推前。若 \(C_{i,r,\omega}\) 表示第 \(i\) 个分支中由前 \(r\) 个二比特对固定为 \(\omega\in(\{0,1\}^2)^r\) 所确定的长度 \(r\) 圆柱集合，则对一切 \(n\ge r\) 都有精确恒等式
\[
\mu_n(C_{i,r,\omega})=\frac{1}{5}\cdot 4^{-r}.
\]
因而圆柱质量自深度 \(r\) 起即严格稳定。进一步，\((\mu_n)_{n\ge 0}\) 在 \(\mathfrak X_{\mathrm{LY}}\) 上弱收敛到概率测度
\[
\mu_\infty
=
\frac{1}{5}\sum_{i=0}^{4} m_i,
\]
其中 \(m_i\) 是分支分量 \(\mathbf P_i+\ZZ_2^2\) 上的 Haar 概率测度；并且同一圆柱集合满足
\[
\mu_\infty(C_{i,r,\omega})=\frac{1}{5}\cdot 4^{-r}.
\]
\end{theorem}

\begin{proof}
定理 \ref{thm:derived-leyang-profinite-rigidity} 已给出
\[
\mathfrak X_{\mathrm{LY}}
\cong
\bigsqcup_{i=0}^{4}\ZZ_2^2,
\]
而推论 \ref{cor:xi-terminal-zm-leyang-branchpoint-tate-torsor-decomposition} 说明每个分支实际上是某个固定提升 \(\mathbf P_i\) 的平移余类，故得到所示分解
\(
\mathfrak X_{\mathrm{LY}}\cong\bigsqcup_{i=0}^{4}(\mathbf P_i+\ZZ_2^2)
\)。

在第 \(n\) 层，每个分支 \(\Pi_{i,n}\) 与 \((\ZZ/2^n\ZZ)^2\) 同构，因而含有 \(4^n\) 个点；总点数为
\[
|\mathcal R_n|=5\cdot 4^n.
\]
把 \((a,b)\in(\ZZ/2^n\ZZ)^2\) 送到其二进展开后接零尾所得到的 \(\ZZ_2^2\) 元素，便得 \(\mathcal R_n\hookrightarrow\mathfrak X_{\mathrm{LY}}\) 的规范嵌入。对任意固定的长度 \(r\) 圆柱集合 \(C_{i,r,\omega}\)，当 \(n\ge r\) 时，在第 \(i\) 个分支内恰有 \(4^{n-r}\) 个点落入该圆柱，于是
\[
\mu_n(C_{i,r,\omega})
=
\frac{4^{n-r}}{5\cdot 4^n}
=
\frac{1}{5}\cdot 4^{-r}.
\]
该值与 \(n\) 无关。另一方面，\(\mathbf P_i+\ZZ_2^2\) 上的 Haar 概率测度正由所有长度 \(r\) 圆柱赋值 \(4^{-r}\) 唯一表征，因此
\[
\mu_\infty(C_{i,r,\omega})
=
\frac{1}{5}\,m_i(C_{i,r,\omega})
=
\frac{1}{5}\cdot 4^{-r}.
\]
由于圆柱集合构成 \(\mathfrak X_{\mathrm{LY}}\) 的紧全不连通拓扑基，并且对每个固定圆柱 \(C_{i,r,\omega}\) 都有
\(
\mu_n(C_{i,r,\omega})=\mu_\infty(C_{i,r,\omega})
\)
在一切 \(n\ge r\) 时成立，故 \(\mu_n\) 在圆柱代数上逐点收敛到 \(\mu_\infty\)，从而 \(\mu_n\Rightarrow\mu_\infty\)。证毕。
\end{proof}
\leanverified{paper\_derived\_leyang\_branchset\_haar\_limit}

\begin{theorem}[Lee--Yang 分支集的 Ihara \texorpdfstring{$\zeta$}{zeta} 函数闭式]\label{thm:derived-leyang-branchset-ihara-zeta}
\leanverified{paper\_derived\_leyang\_branchset\_ihara\_zeta}
把 \(R(y_n)\) 视为有限无向图。则其 Ihara \(\zeta\) 函数满足
\[
Z^{\mathrm{Ih}}_{R(y_n)}(u)
=
\Biggl[
(1-u^2)^{-(2n-2)2^{2n-1}}
\prod_{j=0}^{2n}
\Bigl(1-(2n-2j)u+(2n-1)u^2\Bigr)^{-\binom{2n}{j}}
\Biggr]^5.
\]
\end{theorem}

\begin{proof}
由定理 \ref{thm:derived-leyang-branchset-adjacency-spectrum-heat-trace}，
\[
R(y_n)\cong \bigsqcup_{i=0}^{4}Q_{2n}.
\]
因此只需先计算单个 \(Q_{2n}\) 的 Ihara \(\zeta\) 函数，再取五次方。

对 \(d\)-正则图 \(Q_d\)，标准 Ihara--Bass 行列式公式给出 \cite{TerrasZetaGraphs,StarkTerras1996GraphCoverings}
\[
\bigl(Z^{\mathrm{Ih}}_{Q_d}(u)\bigr)^{-1}
=
(1-u^2)^{|E(Q_d)|-|V(Q_d)|}
\det\!\bigl(I-uA_{Q_d}+(d-1)u^2I\bigr).
\]
这里
\[
|V(Q_d)|=2^d,
\qquad
|E(Q_d)|=d\,2^{d-1},
\]
并且由定理 \ref{thm:derived-leyang-branchset-adjacency-spectrum-heat-trace} 的同一超立方体谱分解，
\[
\operatorname{Spec}(A_{Q_d})
=
\bigl\{
d-2j\ \text{with multiplicity }\binom{d}{j}
:\ 0\le j\le d
\bigr\}.
\]
故
\[
Z^{\mathrm{Ih}}_{Q_d}(u)
=
(1-u^2)^{-(d-2)2^{d-1}}
\prod_{j=0}^{d}
\Bigl(1-(d-2j)u+(d-1)u^2\Bigr)^{-\binom{d}{j}}.
\]
取 \(d=2n\) 即得单个 \(Q_{2n}\) 的闭式。再由 Ihara \(\zeta\) 对不交并取乘法，
\[
Z^{\mathrm{Ih}}_{R(y_n)}(u)=\bigl(Z^{\mathrm{Ih}}_{Q_{2n}}(u)\bigr)^5,
\]
从而得到所示公式。证毕。
\end{proof}

\begin{theorem}[Lee--Yang branch tower 的四元地址、Haar 极限与 Cuntz--Krieger 刚性]\label{thm:derived-leyang-branch-tower-fourary-haar-ck-rigidity}
\leanverified{paper\_derived\_leyang\_branch\_tower\_fourary\_haar\_ck\_rigidity}
记
\[
\mathfrak X_{\mathrm{LY}}:=\varprojlim_n \mathcal R_n,
\]
并沿用定理 \ref{thm:derived-leyang-profinite-rigidity}、推论 \ref{cor:derived-leyang-dyadic-dimension-two}、推论 \ref{cor:derived-leyang-fourary-rigidity} 与定理 \ref{thm:derived-leyang-branchset-haar-limit} 的记号。则在选定相容提升 \(\mathbf P_i\in\Pi_{i,\infty}\) 后，成立如下刚性闭式：
\begin{enumerate}
  \item 逆极限空间分解为
  \[
  \mathfrak X_{\mathrm{LY}}
  \cong
  \bigsqcup_{i=0}^{4}\bigl(\mathbf P_i+\ZZ_2^2\bigr);
  \]
  \item 每个 branch sheet 以及整个 \(\mathfrak X_{\mathrm{LY}}\) 的 Hausdorff 维数与上下 Minkowski 维数都恰等于 \(2\)；
  \item 若 \(\mu_n\) 为有限层的归一化计数测度，则
  \[
  \mu_n\Rightarrow \mu_\infty:=\frac15\sum_{i=0}^{4}m_i,
  \]
  其中 \(m_i\) 为 \(\mathbf P_i+\ZZ_2^2\) 上的 Haar 概率测度；并且对每个长度 \(r\) 的圆柱集合 \(C_{i,r,\omega}\) 都有
  \[
  \mu_\infty(C_{i,r,\omega})=\frac15\cdot 4^{-r};
  \]
  \item 若 \(\sigma\) 表示二比特对地址的左移，则
  \[
  \#\Fix(\sigma^n)=5\cdot 4^n
  \qquad (n\ge 1),
  \]
  从而 Artin--Mazur \(\zeta\) 函数满足
  \[
  \zeta_\sigma(z)
  :=
  \exp\!\left(
  \sum_{n\ge 1}\frac{\#\Fix(\sigma^n)}{n}z^n
  \right)
  =
  (1-4z)^{-5};
  \]
  \item 若定义五份满四元单侧移位的邻接矩阵
  \[
  A_{\mathrm{LY}}:=J_4^{\oplus 5},
  \]
  其中 \(J_4\) 为 \(4\times 4\) 全 \(1\) 矩阵，则其 Bowen--Franks 群与 Cuntz--Krieger 代数分别满足
  \[
  BF(A_{\mathrm{LY}})
  :=
  \coker(I-A_{\mathrm{LY}})
  \cong
  (\ZZ/3\ZZ)^5,
  \qquad
  \mathcal O_{A_{\mathrm{LY}}}
  \cong
  \mathcal O_4^{\oplus 5};
  \]
  \item 任意严格同步且无精确重叠的仿射实现都必须满足地址字母表大小
  \[
  |D|=4.
  \]
\end{enumerate}
\end{theorem}

\begin{proof}
第一条正是定理 \ref{thm:derived-leyang-branchset-haar-limit} 中的逆极限分解。第二条由推论 \ref{cor:derived-leyang-dyadic-dimension-two} 立即给出。第三条则是定理 \ref{thm:derived-leyang-branchset-haar-limit} 的精确内容。

由定理 \ref{thm:derived-leyang-profinite-rigidity}，
\[
\mathfrak X_{\mathrm{LY}}
\cong
\bigsqcup_{i=0}^{4}\bigl(\{0,1\}^{2}\bigr)^{\NN},
\]
并且 \(\sigma\) 在每个分支上都共轭于四符号满移位。故每个分支上 \(\sigma^n\) 的不动点恰有 \(4^n\) 个，五个分支合计即得
\[
\#\Fix(\sigma^n)=5\cdot 4^n.
\]
代入 Artin--Mazur \(\zeta\) 函数定义，
\[
\zeta_\sigma(z)
=
\exp\!\left(
\sum_{n\ge 1}\frac{5\cdot 4^n}{n}z^n
\right)
=
\exp\!\left(
5\sum_{n\ge 1}\frac{(4z)^n}{n}
\right)
=
(1-4z)^{-5}.
\]

再看第五条。单个满四元单侧移位对应的邻接矩阵是 \(J_4\)，故五个不交并分支对应于块对角矩阵 \(A_{\mathrm{LY}}=J_4^{\oplus 5}\)。对单个块，
\[
I-J_4
=
\begin{pmatrix}
0&-1&-1&-1\\
-1&0&-1&-1\\
-1&-1&0&-1\\
-1&-1&-1&0
\end{pmatrix},
\]
经过整行列初等变换可化为 Smith 标准形
\[
\operatorname{diag}(1,1,1,3).
\]
因此
\[
BF(J_4)=\coker(I-J_4)\cong \ZZ/3\ZZ,
\]
从而
\[
BF(A_{\mathrm{LY}})
\cong
BF(J_4)^{\oplus 5}
\cong
(\ZZ/3\ZZ)^5.
\]
另一方面，\(\mathcal O_{J_4}\) 正是四元 Cuntz 代数 \(\mathcal O_4\)；块对角不交并对应 Cuntz--Krieger 代数的有限直和，因此
\[
\mathcal O_{A_{\mathrm{LY}}}
\cong
\mathcal O_{J_4}^{\oplus 5}
\cong
\mathcal O_4^{\oplus 5}.
\]

最后，第六条正是推论 \ref{cor:derived-leyang-fourary-rigidity} 的结论。证毕。
\end{proof}

\subsection{Lee--Yang 振幅常数的六次全虚提升}

在 Lee--Yang 三次数域已经由 \(u=\kappa_*^2\)、\(b=B^2\) 与 \(\lambda_*\) 共同刚性化之后，幅值常数 \(B\) 本身还可继续提升为一个六次全虚数域的原始生成元。

\begin{theorem}[Lee--Yang 振幅常数 \(B\) 的六次极小多项式与全虚签名]\label{thm:derived-leyang-amplitude-degree-signature}
沿用定理 \ref{thm:xi-terminal-zm-kappa-square-cubic-field-s3}、定理 \ref{thm:xi-terminal-zm-stokes-amplitude-square-cubic-rigidity} 与定理 \ref{thm:xi-terminal-zm-stokes-leyang-shared-artin-representation} 的记号，令 \(B\) 为相应的 Lee--Yang 振幅常数，并记
\[
u:=\kappa_*^2,
\qquad
b:=B^2.
\]
则：
\[
[\QQ(B):\QQ]=6,
\]
\[
\operatorname{minpoly}_{\QQ}(B)
=
162X^6+1593X^4+1998X^2+1184,
\]
并且数域 \(\QQ(B)\) 的签名为
\[
(r_1,r_2)=(0,3).
\]
换言之，\(\QQ(B)\) 是一个六次全虚数域。
\leanverified{paper\_derived\_leyang\_amplitude\_degree\_signature}
\end{theorem}

\begin{proof}
由定理 \ref{thm:xi-terminal-zm-stokes-leyang-shared-artin-representation}，
\[
K:=\QQ(b)=\QQ(u)=\QQ(\lambda_*),
\qquad
[K:\QQ]=3.
\]
又由定理 \ref{thm:xi-terminal-zm-stokes-amplitude-square-cubic-rigidity}，\(b\) 满足不可约三次方程
\[
162b^3+1593b^2+1998b+1184=0.
\]
由 Vieta 公式，
\[
N_{K/\QQ}(b)
=
(-1)^3\frac{1184}{162}
=
-\frac{592}{81}.
\]
若 \(b\in K^{\times 2}\)，则其范数 \(N_{K/\QQ}(b)\) 必为 \(\QQ^\times\) 中的平方；但 \(-592/81<0\)，不可能为有理平方，故 \(b\notin K^{\times 2}\)。因此 \(X^2-b\) 在 \(K[X]\) 中不可约，从而
\[
[\QQ(B):K]=2,
\qquad
[\QQ(B):\QQ]=2\cdot 3=6.
\]

再看多项式
\[
g(X):=162X^6+1593X^4+1998X^2+1184.
\]
由于 \(B^2=b\)，直接代入得
\[
g(B)=162b^3+1593b^2+1998b+1184=0.
\]
而 \(g\) 的次数恰为 \(6=[\QQ(B):\QQ]\)，故 \(g\) 即为 \(B\) 在 \(\QQ\) 上的极小多项式。

最后，\(g(X)\) 只含偶次项且所有系数均为正，因此对任意实数 \(x\) 都有 \(g(x)>0\)。故 \(g\) 无实根，\(\QQ(B)\) 没有实嵌入。由
\(
r_1+2r_2=[\QQ(B):\QQ]=6
\)
立得 \((r_1,r_2)=(0,3)\)。证毕。
\end{proof}

\begin{theorem}[Lee--Yang 三次数域不变量的平方根塔不可压缩性]\label{thm:derived-leyang-cubic-invariants-no-quadratic-radicals}
沿用定理 \ref{thm:xi-terminal-zm-stokes-leyang-shared-artin-representation} 与定理 \ref{thm:derived-leyang-amplitude-degree-signature} 的记号，记
\[
K:=\QQ(u)=\QQ(b)=\QQ(\lambda_*),
\qquad
u:=\kappa_*^2,
\qquad
b:=B^2.
\]
则成立：
\begin{enumerate}
  \item \(u,b,\lambda_*\) 都是同一三次数域 \(K\) 的原始元；
  \item \(K/\QQ\) 不是 Galois 扩张，而其伽罗瓦闭包的 Galois 群同构于 \(S_3\)；
  \item \(u,b,\lambda_*\) 均不属于任何由有限次二次扩张组成的塔。特别地，它们都不可能仅由有理运算与有限次嵌套平方根给出精确闭式。
\end{enumerate}
\end{theorem}

\begin{proof}
由定理 \ref{thm:xi-terminal-zm-stokes-leyang-shared-artin-representation}，
\[
\QQ(u)=\QQ(b)=\QQ(\lambda_*)=K,
\qquad
[K:\QQ]=3.
\]
故 \(u,b,\lambda_*\) 各自都生成 \(K\)，从而都是 \(K\) 的原始元。

同一定理同时给出：\(K\) 的伽罗瓦闭包 \(L\) 满足
\[
\Gal(L/\QQ)\cong S_3.
\]
由于 \(S_3\) 非交换，\(|\Gal(L/K)|=2\)，故 \(K/\QQ\) 不是 Galois 扩张。

最后，若 \(u\) 属于某个由有限次二次扩张组成的塔 \(K_r/\QQ\)，则
\[
[K_r:\QQ]=2^r.
\]
而最小多项式次数 \([\QQ(u):\QQ]=3\) 必整除 \([K_r:\QQ]\)，这与 \(3\nmid 2^r\) 矛盾。故 \(u\) 不可能落在任何纯二次塔中。对 \(b\) 与 \(\lambda_*\) 完全同理，因为
\[
[\QQ(b):\QQ]=[\QQ(\lambda_*):\QQ]=3.
\]
因此三者都不可能由有限次嵌套平方根精确表示。证毕。
\end{proof}

\subsection{Lee--Yang 三次模型的共同二次分辨与 Frobenius 奇偶同步}

在 \(u=\kappa_*^2\)、\(b=B^2\) 与 Lee--Yang 分歧三次的共同三次数域已经被识别之后，其判别式平方类与模 \(p\) 分解奇偶型亦可统一压缩到同一条二次角色上。

\leanverified{paper\_derived\_leyang\_cubic\_models\_common\_quadratic\_resolvent}
\begin{theorem}[Lee--Yang 三次模型的共同二次分辨场]\label{thm:derived-leyang-cubic-models-common-quadratic-resolvent}
定义三次多项式
\[
f_u(U):=324U^3-540U^2+333U-74,
\]
\[
f_b(T):=162T^3+1593T^2+1998T+1184,
\]
\[
P_{\mathrm{LY}}(Y):=256Y^3+411Y^2+165Y+32.
\]
令 \(L_u,L_b,L_{\mathrm{LY}}\) 分别为其分裂域。则存在同一个非 Galois 三次数域
\[
K=\QQ(u)=\QQ(b)=\QQ(\beta)
\]
使得
\[
L_u=L_b=L_{\mathrm{LY}}=:L,
\qquad
\Gal(L/\QQ)\cong S_3.
\]
并且三条三次模型的判别式分别满足
\[
\Disc(f_u)=-2^6\cdot 3^9\cdot 37,
\]
\[
\Disc(f_b)=-2^2\cdot 3^9\cdot 11^2\cdot 37\cdot 109^2,
\]
\[
\Disc(P_{\mathrm{LY}})=-3^9\cdot 31^2\cdot 37.
\]
因而它们的平方自由部分同为
\[
-111=-3\cdot 37,
\]
并且共同分裂域 \(L\) 的唯一二次子域为
\[
\QQ(\sqrt{-111}).
\]
\end{theorem}
\begin{proof}
由定理 \ref{thm:xi-terminal-zm-stokes-leyang-shared-artin-representation}，
\[
\QQ(u)=\QQ(b)=\QQ(\beta)=:K,
\qquad
\Gal(L/\QQ)\cong S_3,
\]
其中 \(L\) 为 \(K/\QQ\) 的伽罗瓦闭包。于是三条不可约三次模型的分裂域都等于同一个 \(L\)。

对 \(f_u\) 与 \(f_b\)，推论 \ref{cor:xi-leyang-u-b-square-common-quadratic-resolvent} 已给出
\[
\Disc(f_u)=-2^6\cdot 3^9\cdot 37,
\qquad
\Disc(f_b)=-2^2\cdot 3^9\cdot 11^2\cdot 37\cdot 109^2.
\]
对 \(P_{\mathrm{LY}}\)，定理 \ref{thm:xi-time-part9g-leyang-cubic-discriminant} 给出
\[
\Disc(P_{\mathrm{LY}})=-3^9\cdot 31^2\cdot 37.
\]
因此三者的平方自由部分都等于 \(-111\)。

对任一不可约 \(S_3\)-三次多项式，其分裂域中的唯一二次子域由判别式平方类决定，即
\[
\QQ(\sqrt{\Disc})=\QQ(\sqrt{\Disc_{\mathrm{sf}}}).
\]
故三条三次模型所对应的共同分裂域 \(L\) 的唯一二次子域必为
\[
\QQ(\sqrt{-111}).
\]
证毕。
\end{proof}

\begin{corollary}[Lee--Yang 三次模型模 \texorpdfstring{$p$}{p} 分解型的奇偶同步]\label{cor:derived-leyang-cubic-models-parity-synchronization}
沿用定理 \ref{thm:derived-leyang-cubic-models-common-quadratic-resolvent} 的记号。对任意素数
\[
p\nmid \Disc(f_u)\Disc(f_b)\Disc(P_{\mathrm{LY}})
\]
记
\[
\varepsilon(p):=\left(\frac{-111}{p}\right)\in\{\pm 1\}.
\]
\leanverified{paper\_derived\_leyang\_cubic\_models\_parity\_synchronization}
则公共分裂域 \(L/\QQ\) 上的 Frobenius 元满足
\[
\operatorname{sgn}\bigl(\Frob_p(L/\QQ)\bigr)=\varepsilon(p).
\]
并且 \(f_u,f_b,P_{\mathrm{LY}}\) 在 \(\mathbb F_p\) 上的分解型完全同步：
\begin{enumerate}
  \item 若 \(\varepsilon(p)=-1\)，则三者都恰有一个一次因子与一个不可约二次因子，即分解型均为 \((1)(2)\)；
  \item 若 \(\varepsilon(p)=+1\)，则三者同时处于偶型分解：要么同时为 \((1)(1)(1)\)，要么同时为 \((3)\)。
\end{enumerate}
\end{corollary}
\begin{proof}
由定理 \ref{thm:derived-leyang-cubic-models-common-quadratic-resolvent}，\(L/\QQ\) 的唯一二次子域为
\[
\QQ(\sqrt{-111}).
\]
因此 \(S_3\) 的符号特征正对应于该二次扩张上的二次 Dirichlet 角色，故对任意未分歧素数 \(p\) 都有
\[
\operatorname{sgn}\bigl(\Frob_p(L/\QQ)\bigr)=\left(\frac{-111}{p}\right).
\]

另一方面，对不可约三次的公共分裂域而言，Frobenius 共轭类与良约化下的模 \(p\) 分解型一一对应：恒等元对应 \((1)(1)(1)\)，换位对应 \((1)(2)\)，三循环对应 \((3)\)。其中换位为奇置换，而恒等元与三循环都为偶置换。于是当 \(\varepsilon(p)=-1\) 时，Frobenius 必为换位，三条三次模型都具有分解型 \((1)(2)\)；当 \(\varepsilon(p)=+1\) 时，Frobenius 必为恒等元或三循环，故三者都处于偶型分解。

又因为三条三次模型共享同一个分裂域 \(L\)，在所给的良素条件下它们都由同一 Frobenius 共轭类控制，所以分解型必同步出现。证毕。
\end{proof}

\begin{remark}[三条三次模型上的 joint Chebotarev 退化]\label{rem:derived-leyang-cubic-models-no-fiber-product}
定理 \ref{thm:derived-leyang-cubic-models-common-quadratic-resolvent} 的共同域识别表明：\(f_u,f_b,P_{\mathrm{LY}}\) 并不产生三条彼此独立的 \(S_3\)-闭包，而是同一三次数域 \(K\) 的三个不同三次模型。因而在这三条模型之间不存在非平凡的
\[
S_3\times_{C_2}S_3\times_{C_2}S_3
\]
型 joint Galois 群；相应的联合素数分解律已退化为单一 \(S_3\)-扩张 \(L/\QQ\) 的普通 Chebotarev 分布。
\end{remark}
